% Auto-generated by tables/generate_tables.py
% Source: /home/agagora/Downloads/GITHUB/LLM-HypatiaX-DEV/hypatiax/data/results/comparison_results/feynman-tests/noise-sweep/protocol_core_*.json
% Date:   2026-10-03 21:04
% DO NOT EDIT MANUALLY — re-run 'make tables' to regenerate

\begin{table}[H]
\centering
\caption{Internal routing statistics --- evidence that M3 and M4 are genuinely distinct architectures. Regenerated by \texttt{generate\_tables.py} from the 5 \texttt{protocol\_core} result files ($\sigma\in\{0, 0.05, 0.1, 0.5, 1\}\%$; 30 equations per run, so 150 records per system). ``Ensemble-first'' (M3) and ``LLM-first'' (M4) denote the share of records whose recorded \texttt{metadata.decision} routed to that branch; this is a routing share, not a solve rate, and the \texttt{success} flag is not used. The pool mixes the noiseless run with the noisy runs.}
\label{tab:arch}
\small
\begin{tabular}{L{3.5cm} r r}
\toprule
\textbf{Decision path} & \textbf{M3 (\EHD)} & \textbf{M4 (\HSL)}\\
\midrule
Primary path (routing share) & 140/150 (93.3\%) ensemble-first & 123/150 (82.0\%) LLM-first\\
Alternative path & 10/150 (6.7\%) NN & 27/150 (18.0\%) ensemble\\
Secondary path & NN refinement & NN residual correction\\
Architecture type & Ensemble-first & LLM-first + NN-refinement\\
\bottomrule
\end{tabular}
\end{table}
